\section{Scope and literature decision gate}\label{sec:scope-gate}

\subsection{Exact scope}

\Cref{thm:eventual-selected-return} treats finite-dimensional,
fixed-maximum-delay RFDEs on a continuous-history space.  It deliberately
excludes state-dependent or neutral delays, infinite delays, infinite
networks, and infinite-dimensional node states.  It asserts the strict
inequality \(T_->k\tau_*\) only as a sufficient regularization threshold.
It does not prove that this threshold is necessary or optimal.

The theorem is conditional on a common solution tube, terminal event domain,
strict endpoint signs, and uniform positive speed.  It does not manufacture
those inputs for a concrete model.  Likewise,
\cref{lem:direct-selected-stable-germ} is a local identification result.  It
does not prove hyperbolicity, construct a stable graph, establish a global
basin boundary, or provide recurrent selected hits.

\subsection{What must be checked before any novelty claim}

The books cited below provide standard RFDE semiflow, smooth-dependence, and
Poincar\'e-map background, but this draft has not yet carried out the
theorem-by-theorem comparison required to locate the two statements.  A
future literature audit must answer at least the following questions.

\begin{enumerate}[label=(L\arabic*)]
\item Is joint \(C^k\) dependence of
  \((t,\phi)\mapsto\Phi_t(\phi)\) on the continuous-history phase space after
  \(k\tau_*\) already stated in this form, or is the induction here an
  adaptation of fixed-time regularity results?
\item What are the sharpest established hypotheses for smooth event-defined
  RFDE Poincar\'e maps, especially when the event is selected in a late common
  window rather than required to be the first hit?
\item Is the stable-germ conclusion for a bounded positive-time selected
  return an explicit published lemma, an immediate suspension argument, or a
  useful statement that is usually left implicit?
\item Which common-domain and compatibility hypotheses in the published
  results correspond exactly to the \(U\)-tube, terminal event domain,
  retained local section, and recurrent-hit assumptions used here?
\end{enumerate}

Until these questions are resolved from primary sources, the correct status
is: two isolated working arguments with possible reusable value, but no
established independent originality and no venue claim.

\begin{thebibliography}{9}

\bibitem[Diekmann et al.(1995)]{diekmann1995delay}
O.~Diekmann, S.~A. van Gils, S.~M. Verduyn Lunel, and H.-O. Walther.
\newblock \emph{Delay Equations: Functional-, Complex-, and Nonlinear
  Analysis}.
\newblock Applied Mathematical Sciences 110, Springer, New York, 1995.
\newblock \href{https://doi.org/10.1007/978-1-4612-4206-2}
  {doi:10.1007/978-1-4612-4206-2}.

\bibitem[Hale and Verduyn Lunel(1993)]{hale1993introduction}
J.~K. Hale and S.~M. Verduyn Lunel.
\newblock \emph{Introduction to Functional Differential Equations}.
\newblock Applied Mathematical Sciences 99, Springer, New York, 1993.
\newblock \href{https://doi.org/10.1007/978-1-4612-4342-7}
  {doi:10.1007/978-1-4612-4342-7}.

\end{thebibliography}
